\subsection{Islahat Fermanı}
\indent\indent Her ne kadar Tanzimat'ın uygulanmasında sorunlar yaşansa da, devleti Islahat Fermanı'na götüren süreç Kırım Savaşı ile başladı. Katolikler ve Ortodokslar arasında çıkan anlaşmazlık, Fransa'nın 
Katoliklerin, Rusya'nın da Ortodoksların hamiliğine soyunması durumu uluslararası bir kriz haline getirdi. Osmanlı Devleti'nin iki devleti de tatmin edece bir çözüm bulamaması üzerine Rusya, isteklerini 
güç kullanarak kabul ettirme yolunu seçmiş ve 1853 Haziranı'nda Rus birlikleri Prut Nehri'ni geçmiştir. 30 Kasım 1853'te Rus Donanması'nın Sinop Limanı baskını, İngiltere ve Fransa'nın arabuluculuğuna ön ayak 
olmuştur. İngiltere ve Fransa, Rusya'dan dört noktada garanti talep etti:
\begin{enumerate}
    \item Tuna Prenslikleri üzerindeki koruyuculuk iddiasından vazgeçmesi
    \item Tuna Nehri'nin ticarete açılması
    \item 1841 Londra Boğazlar Sözleşmesi'ne uyulması
    \item Osmanlı Devleti'ndeki Ortodokslar'ın hamiliğini bırakması
\end{enumerate} 

Buradaki dördüncü madde karşılığında, Osmanlı Devleti'den de vatandaşlarına eşit haklar tanımasını ve karma mahkemeler kurması talep ediliyordu. Tanzimat Fermanı her ne kadar bilâ-istisna herkes kanunlar önünde 
eşittir dese de, bu durum uygulamaya yansımamıştı. Gayrimüslimlerin müslümanlara karşı şehâdeti, müslümanlara ait mahkemelerde kabul görmüyordu.\footcite[284]{tanzimat} Öte yandan cizyeyi ödemeye devam ediyorlar ve 
devlet memurluklarında görev alamıyorlardı. Bu da gayrimüslimlerin, müslümanlarla tam bir eşitlik içinde bulunmasını engelliyordu.

Paris Barış Konferası'na giderken, biraz da aceleye gelmiş bir şekilde Islahat Fermanı ilan edildi. Ferman, ana hatlarıyla Tanzimat'ta tarif edilen vatandaşlık haklarını onaylıyor ve eksiklerini tamamlıyordu. Gayrimüslimlere 
askerlik ve memurluk yolu açılması, vatandaşlık kimliğinin kağıt üzerimde tescili olmuştur. Kanun önünde eşitlik, cezaların türleri ve infaz şekillerine de yansıyordu.
Bunun hukuki meali ise zımmîliğin sona ererek bütün teb`ayı kapsayacak bir ceza hukukunun teşekkülü idi.\footfullcite[164]{mapusane}

